
\subsection{Methodik der experimentellen Betrachtung}
Durch die Literaturrecherche konnten bereits verschiedene
Aspekte zur zeitlichen Optimierung von \acp{TEA} identifiziert
werden. Vorrangig
wurde dafür der verwendete Covert-Channel - meist \ac{CbCC} -
identifiziert. Um die Möglichkeiten zur Optimierung und
damit die Grenzen von \ac{CbCC} im Kontext von Transient
Execution zu identifizieren, wurde ein Test-Programm entwickelt,
welches die zwei verschiedenen Encodierungen mit verschiedenen
Parametern misst. Auf dieser Basis werden die unterschiedlichen
Techniken verglichen und die Parameter ermittelt, welche
hardwarespezifisch betrachtet werden müssen.
Zeitliche Messungen werden mit der
\verb|rdtsc|-Instruktion aufgenommen und deshalb in
\ac{CPU}-Zyklen angegeben.

\paragraph{Der Programmablauf} des Test-Programms für $i$
Trainingsläufe und $n$ übertragene Bits ist im Folgenden
dargestellt:

\begin{figure}[H]
    \begin{enumerate}
        \item Ermittlung zeitliche Grenze gecacht/nicht gecacht (tiefer in
        Anhang \ref{appendix:explanation-on-cache-measurements} erläutert)
        \item \tikzmark{start-measurement}Allokation des Trainingsdaten- ($i+1$ Elemente) und
        Covert-Channel-Arrays
        \item \label{enumerate:set-training-values}Trainings-Array
        Wert für Index $j$ auf Wert $j$ setzen
        \item Flushen des Training- und Covert-Channel-Arrays
        (Trainings-Array wird nur Index $i$ geflusht) und Laden des Secrets
        \item Trainieren des Predictors (Ausführen von Transfer mit Index
        $0$ bis $i-1$)
        \item \label{enumerate:transient-function}Prüfung, ob
        $Array[i]<i$ führt zu spekulativer Ausführung
        \begin{itemize}
            \item \tikzmark{start-transient}Schleife über Anzahl übertragener
            Bits \textrightarrow\space Aufbau Bitmaske
            \item Für bitweise Encodierung jede Iteration laden, bzw. nicht
            laden des Index
            \item Für Array-Index-Encodierung, Bitmaske über Secret legen
            \textrightarrow\space Index aus Covert-Channel-Array laden
            (einmalig nach Aufbau der Bitmaske)
            \item $Array[i]<i$ wird evaluiert, sobald $Array[i]$ geladen
            \textrightarrow\space führt zu Rollback (da Wert
            \tikzmark{end-transient}entsprechend
            Schritt \ref{enumerate:set-training-values} gesetzt)
        \end{itemize}
        \item \tikzmark{end-measurement} Auslesen des Secrets aus Cache-Zustand
        \item Schreiben der Werte in die Ergebnisdatei
    \end{enumerate}
    \begin{tikzpicture}[remember picture,overlay]
        \coordinate (a) at ({pic cs:start-measurement});
        \coordinate (b) at ({pic cs:end-measurement});
        \coordinate (c) at ({pic cs:start-transient});
        \coordinate (d) at ({pic cs:end-transient});
        \draw [decoration={brace,mirror},thick,decorate] ([yshift=2ex,xshift=2ex]a -| current page text area.west) -- node [midway,left,rotate=90,anchor=south,inner sep=5pt] {zeitliche Messung}  ([yshift=-.5ex,xshift=2ex]b -| current page text area.west);
        \draw [decoration={brace,mirror},thick,decorate] ([yshift=2ex,xshift=7ex]c -| current page text area.west) -- node [midway,left,rotate=90,anchor=south,inner sep=5pt] {Transient Execution}  ([yshift=-.5ex,xshift=7ex]d -| current page text area.west);
    \end{tikzpicture}
    \caption{Aufbau des Test-Programms\label{figure:test-programm-steps}}
\end{figure}

\paragraph{Als Encodierung in den \acs{CbCC}}
verwenden die meisten \acp{TEA}
die Array-Index-En\-co\-die\-rung. (\cite{lipp_meltdown_2020},
\cite{van_schaik_ridl_2019}, \cite{schwarz_zombieload_2019},
\cite{hertogh_rain_nodate}, \cite{kim_slap_2025},
\cite{kim_flop_2025})
Grund dafür ist zum einen, dass
mehr als ein Informationsbit mit dem Laden einer einzelnen
Adresse übertragen werden kann, und zum anderen, dass die
ursprüngliche Spectre-Beschreibung häufige Muster in Software
als Gadgets ausnutzt, bei denen solch ein Speicherzugriff von
dem Secret abhängig ist (\cite[8]{kocher_spectre_2019}).

Obwohl bei diesen Angriffen das Gadget aus Code besteht, der vom
Angreifer kontrolliert wird, sind mir bisher keine praktischen
Versuche in der Literatur bekannt, das Secret binär, mit mehr als
einem Bit, zu encodieren. Um die Grenzen von \ac{CbCC}
abzuschätzen, müssten potenzielle Wege, den Durchsatz dieser
Kanäle zu erhöhen, gefunden und evaluiert werden. Aus diesem Grund
ist ein Ziel dieser Arbeit, Daten binär in einen \ac{CbCC} während
Transient Execution zu encodieren, um die praktische
Umsetzbarkeit zu prüfen und verschiedene technische
Eigenschaften zu evaluieren und zu vergleichen.

\paragraph{Als Covert-Channel-Technik}
wurde in dem Test-Programm der Flush+Reload Covert-Channel
(\cite{yarom_flushreload_2014}) verwendet,
da das Ziel primär das Demonstrieren von binärer Encodierung in
einem evictionbasierten \ac{CbCC}, nicht die Wahl des
Covert-Channels mit dem höchstmöglichen Durchsatz ist. Andere
Covert-Channel-Techniken können ein anderes Laufzeitverhalten,
sowie andere Fehlerraten und mehr verfügbare Bits aufweisen.
Genauso kann für einige Covert-Channel die binäre Encodierung
einen größeren oder kleineren Vorteil bieten als für
Flush+Reload. Beispielsweise kann die Instruktion für das
Encodieren eines Bits eine längere oder kürzere Laufzeit haben
und damit mehr oder weniger Bits in einem transient Fenster
übertragbar sein.

\paragraph{Sicherstellung von Transient Execution.}
Um Transient Execution auszulösen, wurde eine Branch-Instruktion
verwendet, die außerdem, falls nötig, einen Predictor
trainieren kann. Die Funktion, in der sich das transient
Fenster befindet, ist in Listing \ref{listing:transient-function}
zu sehen. In Abbildung \ref{figure:test-programm-steps} findet
dies in Schritt \ref{enumerate:transient-function} statt. Dort
wird abhängig von einem Wert aus einem Array
geprüft, ob Daten in den \ac{CbCC} geschrieben werden sollen
oder nicht. Das Array hält für jeden Index $i$ (beginnend mit 0)
den Wert $i+1$, sodass, für alle, außer den letzten Zugriff,
die Branch-Bedingung gilt. Das letzte Element der Trainingsdaten wird
zudem vor dem Training mit der \verb|clflush|-Instruktion aus dem
Cache entfernt, damit die \ac{CPU} entweder den korrekten Pfad
vorhersagen, oder die Ausführung verzögern muss.
Während der Trainingsphase, insofern
eine angegeben ist, wird als Secret immer 0 in den \ac{CbCC}
geschrieben, damit keine Adressen geladen werden. Wenn von der
Methode das Encodieren in den \ac{CbCC} aufgerufen wird, wird
der boolesche Wert „true“ als Indikator für den Trainingsprozess
zurückgegeben. In der Übertragungsphase wird „false“
zurückgegeben, um sicherzugehen, dass aus architektonischer Sicht
\verb|asm_encode_data| nicht ausgeführt wurde. Wird aus dem
\ac{CbCC} dennoch das Secret gelesen, wurden die Instruktionen
als Transient Execution ausgeführt.

\paragraph{Encodierung in den Cache-Zustand.}
Listing \ref{listing:cbcc-binary-encoding} zeigt die Encodierung
des Secrets in binärer Form in den Cache-Zustand. In Abbildung
\ref{figure:test-programm-steps} findet dies bei den Punkten
6 und 7 statt. Dabei
wird ein Array als Eingabe genommen, dessen Elemente vorher durch
\verb|clflush| aus dem Cache entfernt wurden. In diesem Array
wird mit einem Abstand von \verb|stride| eine Adresse geladen,
wenn diese den Wert 1 annehmen soll, ansonsten wird diese
übersprungen. Zu beachten sei hier, dass die erste
Datenadresse mit einem Abstand von \verb|stride| zum Beginn des
Arrays gewählt wurde, um Beeinflussungen des Cache-Zustands durch
vorherige Objekte auf dem Stack zu umgehen. Es kann außerdem
angegeben werden, wie viele Bits encodiert werden sollen.

Bei der Array-Index-Encodierung wird genauso ein Array übergeben,
allerdings wird der Datenwert, nach dem Verrechnen mit einer
Bitmaske, mit dem \verb|stride| multipliziert und das Ergebnis als
Index für die zu ladende Adresse in dem Array genommen.

\paragraph{Ausgeführt} wurde das Programm auf verschiedenen
\acp{CPU}, von denen die Messwerte in Anhang
\ref{appendix:measurements-graphs} zusammengefasst werden.

Ein auffälliger Aspekt ist, dass der \ac{CbCC} bei einem
AMD Ryzen 7 7735HS keine oder nur sehr geringe
Daten korrekt transferiert hat. Die Genauigkeit befand sich
dauerhaft knapp unter 0,6. Während das Problem die
Implementierung des Programms sein könnte, weist der
Covert-Channel auf der AMD A10-7850K APU deutlich höhere
Genauigkeiten auf. Dies wirft die Frage auf, ob neuere \acp{CPU}, wie
der Ryzen 7 7735HS, Mitigationen gegen \ac{CbCC} aufweisen. Da
mir keine Quelle bekannt ist, die solch
eine Mitigation bei diesen \acp{CPU} beschreibt, scheint dies eine
offene Forschungsfrage zu sein. Es stellt sich die Frage, ob
neuere Ryzen-\acp{CPU} Mitigationen gegen \ac{CbCC},
möglicherweise nur während Transient Execution, besitzen, und
wenn dies der Fall ist, welche Arten von Angriffen dadurch
verhindert werden.
Eine vollständige Mitigation von \acp{TEA} durch das Verhindern
von Covert-Channels ist zu bezweifeln, da contentionbasierte
Covert-Channels nahezu unvermeidbar sind und die einzige bekannte
garantierte Mitigation von \acp{TEA} das
Deaktivieren von Transient Execution ist. Die daraus
resultierenden Performanceeinbußen wären zu groß, um die
Verhinderung dieser Angriffe zu rechtfertigen.

\subsection{Auswertung}
\label{section:experimental-results}
Die Abbildungen
\ref{figure:measurement-results-bitwise},
\ref{figure:measurement-results-array-index},
\ref{figure:measurement-runtime-bitwise} und
\ref{figure:measurement-runtime-array-index}
zeigen die Ergebnisse als Graphen zusammengefasst.

\paragraph{In diesen Grafiken}
bedeutet „bitweise korrekt“, dass das Verhältnis
aller korrekt übertragenen Bits, über alle Übertragungen hinweg
gegenüber der Anzahl an übertragenen Bits ins Verhältnis gesetzt
wird. Gesamte Übertragungen korrekt setzt dagegen die Anzahl
an Übertragungen, bei denen kein Fehler aufgetreten ist,
gegenüber der Anzahl an Übertragungen ins Verhältnis. Nötig ist
dies, da bei der bitweisen Encodierung
zwar seltener alle Bits korrekt übertragen werden, allerdings eine
einzelne Änderung des Cache-Zustands, nicht alle Informationen
stört. Bei der Array-Index-Encodierung dagegen wird bei
jeglicher Störung keine Information korrekt übertragen. Die
Fehlerkorrektur bei der bitweisen Encodierung, durch
Wiederholung der Übertragung ist also generell besser möglich.
Bei den beiden Darstellungen sei zu beachten, dass bei keinem
Signal, das bedeutet bei rein zufällig erhaltenen Werten, die
Genauigkeit bitweise betrachtet gegen 0,5 geht, während bei
der Betrachtung der korrekten Übertragungen die Genauigkeit
gegen 0 geht. Die Skalen sind also nicht identisch zu
interpretieren.

Die Abbildungen \ref{figure:measurement-runtime-bitwise} und
\ref{figure:measurement-runtime-array-index} zeigen außerdem
die durchschnittliche Laufzeit einer Übertragung, gemessen
über den in Abbildung \ref{figure:test-programm-steps}
gezeigten Abschnitt.

\paragraph{Die Fehlerwahrscheinlichkeit}\label{paragraph:error-probability}
zwischen den beiden Encodierungen verhält sich ähnlich in einem
Bereich von 0,8 bis 0,95 über alle Trainingslängen hinweg.
Die Array-Index-Encodierung weist grundsätzlich eine höhere
Genauigkeit auf, wobei die Genauigkeit der einzelnen Bits nahezu
identisch mit der Genauigkeit der Übertragungen übereinstimmt
und nicht von der Anzahl übertragener Bits, sondern nur von dem
Stride abhängt. Bei der bitweisen Encodierung zeigt sich dagegen,
besonders bei dem i7-6700K, dass
die Genauigkeit ab 9 Bit der vollständig korrekten
Übertragungen, deutlich abfällt. Dies weist darauf hin, dass bei
diesem Prozessor, die maximale Anzahl an
\verb|prefetch0|-Instruktionen, die während der Transient Execution
ausgeführt werden, 8 beträgt.

Bei der bitweisen Encodierung ist ein großer Vorteil, dass mehr
unabhängige Datenpunkte übertragen werden, als bei der
Array-Index-Encodierung. Dadurch kann der häufigste Wert für jedes
Bit einzeln, anstatt für den Gesamtwert, ermittelt werden. Dadurch
könnte bei schlechterer Genauigkeit, verglichen mit der
Array-Index-Encodierung, mit derselben Anzahl an Wiederholungen der
korrekte Wert übertragen werden. Dafür müssen allerdings alle Bits
gleich wahrscheinlich einen Fehler aufweisen. Im Folgenden wird davon
ausgegangen, dass Fehler bei allen Bits gleich wahrscheinlich
sind.

Verschiedene Aspekte der Messungen zeigen, dass eine
Generalisierung der Fehlerwahrscheinlichkeit abhängig von
gegebenen Parametern schwer ist. Die Ergebnisse der
AMD A10-7850K APU zeigen bei der bitweisen Encodierung für
ein bis drei encodierte Bit deutlich geringere Genauigkeiten
als für höhere Bitzahlen. Bei der Array-Index-Encodierung
dagegen ist bei dieser \ac{CPU} für ein bzw. zwei encodierte Bits
die höchste Genauigkeit aller Übertragungen verzeichnet.
Auffällig ist dabei auch, dass ab einem Stride von $2^{10}$ die
Genauigkeit nur noch sinkt. Während der beobachtete Unterschied
zwischen der bitweisen und der Array-Index-Encodierung von dem
Unterschied in der Encodierung abhängt, wird in dieser Arbeit
kein Versuch unternommen, diese beobachteten Phänomene zu
erklären. Sie zeigen aber, welche Problematik die Betrachtung
der Fehlerwahrscheinlichkeit für die Ermittlung des Durchsatzes
darstellt.

\paragraph{Der Stride}
hat wie erwartet einen deutlichen Einfluss auf die Genauigkeit.
Interessant ist jedoch, dass bei dem Intel\textregistered\space
Core\texttrademark\space i7-6700K ab einem Stride von $2^{9}$ kaum
eine Verbesserung der Genauigkeit zu erkennen ist. Auch bei den
anderen \acp{CPU} zeigt sich ein ähnliches Verhalten, was auf die
maximale Reichweite des Prefetchers hindeutet. Bei einem Stride
von $2^{12}$, was der Größe einer Memory-Page entspricht, lassen
sich jedoch zuverlässig die besten Ergebnisse beobachten.

\paragraph{Die Anzahl der Trainingsläufe}
hat, wie in Abbildung \ref{figure:measurement-results-array-index}
zu sehen, bei der Array-Index-Encodierung keinen relevanten
Einfluss auf die Genauigkeit der Übertragung. Bei der bitweisen
Encodierung dagegen ist eine leicht steigende Tendenz der
Genauigkeit bei steigender Anzahl an Trainingsläufen zu erkennen.
Eine mögliche Erklärung dafür wäre, dass mit mehr Trainingsläufen
die Länge des transient Fensters steigt und damit auch die Anzahl
an korrekt übertragenen Bits bei der bitweisen Encodierung. Bei
der Array-Index-Encodierung dagegen ändert die Größe des
transient Fensters nichts mehr, sobald dessen Dauer für die
Encodierung einer Adresse ausreicht.

\paragraph{Das Laufzeitverhalten}
ist als logarithmische Skala dargestellt, da die Laufzeit der
Ar\-ray-Index-Encodierung mit zunehmender Bitanzahl exponentiell
skaliert. Grund dafür ist, dass, wie bereits erwähnt, die Anzahl
der zu messenden Adressen exponentiell mit der Anzahl an
übertragenen Bit steigt.

Die Abbildungen
\ref{figure:measurement-runtime-bitwise} und
\ref{figure:measurement-runtime-array-index} zeigen, dass der Stride
einen Einfluss auf die Laufzeit der Übertragung hat. Die Ursache
dafür ist, dass in dem Test-Programm die Allokation
von dem Covert-Channel-Array sowie das Setzen der Array-Elemente auf einen
festen Wert mitgemessen werden und deren Laufzeit
von der Größe des Arrays abhängt. Das Setzen von Werten
in dem Array ist nötig, da ansonsten bei einem Stride von 4096 die
Genauigkeit deutlich sinkt. Der Grund dafür könnte mit der Handhabung
von Memory-Pages zusammenhängen, da sich bei diesem Stride nur eine
betrachtete Adresse je Page befindet. In einem Angriff
mit mehreren Übertragungen kann dasselbe Array mehrfach verwendet
werden. Aus diesem Grund ist die Laufzeit dieser Operationen nicht für
das Abschätzen einer Grenze für die Laufzeit von \acp{TEA} relevant.
Für die bitweise Encodierung zeigt sich deutlich, dass der
Covert-Channel kaum einen Einfluss auf die Laufzeit einer Übertragung
hat. Dies zeigt sich, da nur ein geringer Anstieg der Laufzeit einer
Übertragung bei mehr Bit erfolgt. Bei der Array-Index-Encodierung
dagegen steigt die Gesamtlaufzeit mit steigender Bitanzahl deutlich
an. Dies deutet bereits darauf hin, dass die optimale Encodierung
eine Kombination der beiden ist.

\subsection{Theoretische Laufzeitbetrachtung}
Im Folgenden wird die Laufzeit einer Übertragung während eines
\ac{TEA} unter einigen Annahmen und ohne die Genauigkeit des
\ac{CbCC} zu beachten betrachtet. Dabei wird als Zeiteinheit
immer implizit \ac{CPU}-Zyklen verwendet.

\paragraph{Folgende Annahmen werden getroffen:}
\begin{itemize}
    \item Die ausgeführten Operationen dauern bei gleichem
    Zustand der Mikroarchitektur immer exakt gleich
    lang, gemessen an den \ac{CPU}-Zyklen.
    \item Alle möglichen übertragenen Informationen
    sind gleich wahrscheinlich.
    \item Der Unterschied der Dauer des Auslesens einer Adresse,
    wenn diese gecacht ist, gegenüber der Laufzeit, wenn diese
    nicht gecacht ist, wird nicht unterschieden. Für die
    bitweise Encodierung kann zwar, unter der Annahme, dass ein
    Informationsbit mit einer Wahrscheinlichkeit von
    $\frac{1}{2}$ den Wert $1$ hat, der Mittelwert der beiden
    verwendet werden, allerdings gilt dies nicht für die
    Array-Index-Encodierung. Als Zugriffszeit, nachfolgend $T_R$,
    wird die kürzere der beiden Zeiten betrachtet, um eine obere
    Grenze für den Durchsatz durch Überschätzen zu ermitteln.
    \item Eine Übertragung kann beliebig oft wiederholt werden,
    um das Secret schrittweise zu übertragen.
    \item Nach dem Auslesen der letzten Übertragung
    und der nächsten Setup-Phase vergeht keine Zeit.
\end{itemize}

\paragraph{Die Setup-Phase}
muss in der Regel eine Operation, wie das Evicten von
Adressen für alle beobachteten Adressen durchführen.
Zusätzlich muss der Predictor, sofern einer verwendet wird,
trainiert werden. Damit ergibt sich die Laufzeit $L_S$
mit $N_T$ als Anzahl der Trainingsläufe, $T_T$ als die
Dauer eines Trainingslaufs, $N_P$ als die Anzahl der beobachteten
Adressen und $T_P$ als Dauer der Vorbereitung einer Adresse.
\[L_S=N_T*T_T+N_P*T_P\]
Die Vorbereitung einer Adresse besteht beispielsweise bei
Flush+Reload in der Dauer einer Flush-Instruktion.

\paragraph{Für die Transient-Execution-Phase}
wird der zeitliche Aufwand $L_E$ als konstanter Wert betrachtet,
da der Wert vorrangig von der Kondition abhängt, die das transient
Fenster auslöst.

\paragraph{In der Decode-Phase}
hängt die Laufzeit primär von dem verwendeten
Covert-Channel ab. Diese Laufzeit $L_D$ hängt von der Anzahl
auszulesender Adressen $N_R$ und dem zeitlichen Aufwand zum
Auslesen einer Adresse $T_R$ ab.
\[L_D=N_R*T_R\]
Bei der Setup- und der Decode-Phase wird bewusst zwischen $N_P$ und
$N_R$ unterschieden, da die Anzahl der vorbereiteten Adressen
zwangsläufig alle möglichen Zustandsadressen umfasst. Bei dem
Auslesen kann allerdings für die Array-Index-Encodierung bereits
die erste Adresse, die als gecacht ermittelt
wird, auch als encodierter Wert betrachtet werden.
Um die Berechnung der Anzahl der zu messenden Adressen zu
vereinfachen, werden alle Werte, die in den \ac{CbCC} geschrieben
werden könnten, als gleich wahrscheinlich angenommen.
Daraus folgt, dass im Mittel nach der Hälfte der Adressen
der encodierte Wert gefunden wird.
Für die bitweise Encodierung ändert sich dadurch nichts, da dort
jede Adresse einen Informationswert speichert und somit stets alle
Adressen ausgelesen werden müssen.

\paragraph{Die Anzahl an verwendeten Adressen}
wird für eine verallgemeinerte Kombination aus bitweiser und
Array-Index-Encodierung betrachtet.
Es werden zur Encodierung Ad\-ressen aus $n_b$ Arrays geladen.
Jedes Array encodiert $2^{n_a}$ verschiedene Zustände. Der Wert
$0$ wird encodiert, indem keine Adresse geladen ist. Daraus
ergeben sich für jedes Array $2^{n_a}-1$ Adressen. $n_a$ stellt
dabei die Anzahl an Array-Index encodierten Adressen dar.
Damit ergeben sich für $N_P$ und $N_R$ folgende Berechnungen:
\[N_P=n_b*(2^{n_a}-1)\]
\[N_R=n_b*2^{n_a-1}\]
Wobei $n_b$ und $n_a$ mindestens $1$ sind. Die vollständige
bitweise Encodierung liegt also vor, wenn $n_a=1$ und $n_b$
die Anzahl an übertragenen Bits angibt. Wenn $n_b=1$ und
$n_a$ beliebig, liegt die Array-Index-Encodierung vor.
Damit wird davon ausgegangen, dass bei der
Array-Index-Encodierung der Wert 0 übertragen wird, indem
keine andere Adresse gecacht ist.

\paragraph{Die theoretische Laufzeit}
$L_G$ einer einzelnen Übertragung mit den benannten Faktoren
beträgt somit:
\[L_G=L_S+L_E+L_D\]
\[L_G=N_T*T_T+n_b*(2^{n_a}-1)*T_P+L_E+n_b*2^{n_a-1}*T_R\]
Da für einen spezifischen \ac{TEA} die Transient Execution
und deren Vorbereitung gleichbleiben, gilt für
hinreichend große $n_b$, $n_a$, $T_P$ und $T_R$:
\[L_G\propto n_b*(2^{n_a}-1)*T_P+n_b*2^{n_a-1}*T_R\]
Begrenzt ist $n_b$ durch die Tiefe des transient Fensters und
$n_a$ durch die Größe des Cache, beziehungsweise des Speichers.

Der Durchsatz $D$, welcher optimiert werden soll, ergibt sich aus
der Anzahl übertragener Bits $n$ je Zeiteinheit.
\[D=\frac{n}{L_G}=\frac{n_a*n_b}{L_G}\]
Da $L_G$ durch $n_b$ linear und durch $n_a$ exponentiell skaliert,
sollte $n_b$ maximal und $n_a$ minimal gesetzt werden, wobei $n_b$
durch das transient Fenster begrenzt ist. Unter der theoretischen
Annahme, dass $n_b$ nicht begrenzt ist, ergibt sich, da $N_T*T_T$
und $L_E$ konstant sind, für $D$ bei $n_a=1$ die folgende
theoretische Grenze.
\[\lim_{n_b\to +\infty}D=\lim_{n_b\to +\infty}\frac{n_b}{L_G}=
\lim_{n_b\to +\infty}\frac{n_b}{N_T*T_T+n_b*T_P+L_E+n_b*T_R}=
\frac{1}{T_P+T_R}\]
Diese Grenze ergibt sich auch trivial durch die Überlegung, welche
Operationen bei der bitweisen Encodierung je Bit ausgeführt werden
müssen. Unter der Vernachlässigung des zeitlichen Aufwands für das
Verändern des Cache-Zustands während der Transient Execution ergibt
sich, dass der Zeitaufwand für die Übertragung eines Bits die Summe
aus der Vorbereitung und dem Auslesen ist. Folglich ergibt sich der
Durchsatz $\frac{1}{T_P+T_R}$.

Da allerdings $n_b$ begrenzt ist, $N_T*T_T$ und $L_E$ nicht
beide $0$ sind, ergibt sich für einen durch das transient Fenster
begrenzten Wert $n_{bmax}$ die folgende Funktion für $D$.
\[D(n_a)=\frac{n_{bmax}*n_a}{N_T*T_T+L_E+n_{bmax}*T_P*(2^{n_a}-1)+n_{bmax}*T_R*2^{n_a-1}}\]
Als Visualisierung ist für $N_T=50$, $T_T=20$, $L_E=100$,
$T_P=50$, $T_R=120$ und $n_b=8$ die Funktion $D(n_a)$ in Abbildung
\ref{figure:theoretical-runtime} dargestellt.

Das Maximum dieser Funktion ergibt sich aus dem Punkt, an dem
für die Ableitung $D'(n_a)=0$ gilt. Die Ableitung ist für
Vollständigkeit in Abbildung \ref{figure:throughput-derivation}
dargestellt.

\subsection{Diskussion}
Die Wahl des Strides ist für die Datenrate des \ac{CbCC} eher
weniger relevant. Grund dafür ist, dass ein größerer Stride
keine direkten negativen Auswirkungen auf die Datenrate oder
Genauigkeit hat und nur durch die Größe des \ac{RAM} begrenzt
ist. Relevant werden könnte dieser Parameter, wenn das
allokierte Array so groß ist, dass zu viele beobachtete
Adressen in einem Eviction-Set liegen. Beispielsweise, wenn
auch der \ac{LLC} nur wenige Einträge je Eviction-Set aufweist,
die Array-Index-Encodierung verwendet wird und eine Art
invertierte Encodierung verwendet wird. Also, wenn alle verwendeten
Adressen in der Setup-Phase geladen und während der Transient
Execution zur Encodierung einzelne Adressen evicted werden.

Da $n_b$ maximiert werden soll, wäre es sinnvoll, die Laufzeit
der Instruktion zum Verändern des Cache-Zustands während der
Transient Execution zu minimieren und die Encodierung mit
möglichst wenigen Instruktionen zu realisieren. Dadurch könnten
möglichst viele Bits encodiert werden.
Besonders die Genauigkeit der Zeitmessung beeinflusst die
optimale Encodierung. Denn ansonsten muss der Zeitunterschied,
wie beispielsweise in \cite[17]{kim_slap_2025}, verstärkt werden.
Diese Verstärkung erhöht allerdings auch die Laufzeit einer
Messung, da der Unterschied groß genug sein muss, um von dem
ungenauen Zeitgeber unterschieden zu werden.

Für viele Angriffe, die darauf abzielen, Zugriffsbeschränkungen,
beispielsweise einer \ac{VM} oder des Kernels, zu umgehen, ergibt
es deshalb Sinn, eine Kombination aus bitweiser und Array-Index-Encodierung
zu verwenden. Dort ist häufig eine hochauflösende Zeitmessung,
wie zum Beispiel durch die \verb|rdtsc|-Instruktion, möglich.
Dadurch bleibt die Laufzeit des Auslesens relativ gering im
Verhältnis zu der restlichen Übertragung. Beispiele dafür sind
\cite{hertogh_rain_nodate}, \cite{lipp_meltdown_2020},
\cite{maisuradze_ret2spec_2018}.

Bei gewissen Angriffen sind allerdings keine hochauflösenden
Zeitmessungen möglich. Dies ist beispielsweise in Apple Webkit
der Fall (\cite[11]{kim_slap_2025}). Zum Auslesen der Daten aus
dem Covert-Channel wurde von \cite{kim_slap_2025} der
Zeitunterschied zwischen einer gecachten und einer nicht
gecachten Adresse durch logische Gatter verstärkt. Das Auslesen
einer Adresse dauert im Mittel deshalb zwischen 5 und 12
Millisekunden (\cite[17]{kim_slap_2025}). Dadurch ist $T_R$ im
Verhältnis zum restlichen Angriff so groß, dass die Dauer einer
Übertragung nahezu der Laufzeit des Auslesens entspricht.
Das Verwenden der bitweisen Encodierung verringert die Laufzeit
damit deutlich gegenüber der Array-Index-Encodierung. Mit der
Annahme, dass mit jeder Übertragung 8 Bit encodiert werden,
ergibt sich das folgende Verhältnis zwischen den bitweise
encodierten Adressen $N_{Rb}$ und den Array-Index encodierten
Adressen $N_{Ra}$:
\[\frac{N_{Ra}}{N_{Rb}}=\frac{2^{n_a-1}}{n_b}=\frac{128}{8}=16\]
Da das Auslesen praktisch die Laufzeit der gesamten Übertragung
darstellt, wäre in diesem Falle der gesamte Angriff um den Faktor
16 schneller. Dies ist allerdings nur ein theoretischer Wert
und gilt nur für Angriffe, bei denen das Auslesen eine sehr
hohe Laufzeit im Verhältnis zu der restlichen Übertragung
aufweist. Aus Listing 3 und 4 in \cite[6,8]{kim_slap_2025} lässt
sich ablesen, dass ein Byte als Array-Index encodiert wurde.
Die Laufzeit dieses spezifischen Angriffs könnte also unter den
getroffenen Annahmen deutlich optimiert werden und ein
entsprechend höheres Risiko darstellen. Der dort genannte
Durchsatz von 0,384 Bits pro Sekunde (\cite[12]{kim_slap_2025})
könnte theoretisch auf 6,144 erhöht werden.

Das Risiko, das von Angriffen ausgeht, verändert sich mit der
Geschwindigkeit, in der sie durchgeführt werden können. Auch
wenn der Durchsatz deutlich erhöht und damit die Laufzeit des
zuvor beschriebenen \ac{POC} deutlich verringert werden kann,
bleiben \acp{TEA} jedoch aufwändige Angriffe mit hoher Komplexität.
Die Optimierung der verwendeten \acp{CbCC} trägt zwar dazu bei,
die praktische Umsetzbarkeit solcher Angriffe zu erhöhen,
allerdings ist für die Einschätzung nach wie vor vorrangig
die ausgenutzte Schwachstelle und ihre Auswirkungen auf die
Schutzziele relevant. Da aber eine generelle Mitigation für
\acp{TEA} nicht existiert, bleiben diese Angriffe zukünftig
ein Risiko, welches auch von der Laufzeit abhängt. Die ermittelte
Grenze für die Laufzeit eines Angriffs kann bei der Einschätzung
von Angriffen hinsichtlich der praktischen Ausnutzung helfen.
Zusätzlich können Maßnahmen, mit dem Ziel, Covert-Channels während
\acp{TEA} zu begrenzen, hinsichtlich ihrer Wirkung gemessen und
verglichen werden.
